\section{The complete ResNet-20 and its shape flow}
\label{sec:arch}

\subsection{The $6n+2$ family}

\begin{idea}[counting layers]
The CIFAR ResNets of the paper consist of a $3{\times}3$ stem convolution, then three \emph{stages} of $n$ basic blocks each, at feature-map sizes $32$, $16$, $8$ with $16$, $32$, $64$ filters, then global average pooling and a 10-way fully connected layer. Each block holds 2 weighted layers, so
\[
  \underbrace{1}_{\text{stem}} + \underbrace{3\cdot n\cdot 2}_{\text{blocks}} + \underbrace{1}_{\text{FC}} = 6n+2 \text{ weighted layers.}
\]
$n=3$ gives 20 and $n=9$ gives 56. BN, ReLU and pooling are not counted as layers. Only layers with weights are.
\end{idea}

\begin{repo}[\texttt{ResNetCIFAR(n, residual=True, num\_classes=10, in\_channels=3, widths=(16,32,64))}]
\begin{lstlisting}[firstnumber=45]
        self.stem = nn.Conv2d(bias=False,stride=1,padding=1,kernel_size=3,in_channels=in_channels,out_channels=self.widths[0])
        self.bn1 = nn.BatchNorm2d(self.widths[0])
        self.relu = nn.ReLU()
        self.linear = nn.Linear(self.widths[2],self.num_classes)
        self.avg_pool = nn.AdaptiveAvgPool2d(1)
        self.flatten = nn.Flatten(1,-1)
        self.stage1 = self._make_stage(in_ch=widths[0],out_ch=widths[0],stride=1)
        self.stage2 = self._make_stage(in_ch=widths[0],out_ch=widths[1],stride=2)
        self.stage3 = self._make_stage(in_ch=widths[1],out_ch=widths[2],stride=2)
        self._init_weights()
    def _make_stage(self,in_ch,out_ch,stride):
        blocks = []
        blocks.append(BasicBlock(in_ch=in_ch,out_ch=out_ch,stride=stride,residual=self.residual))
        for i in range(self.n - 1):
            blocks.append(BasicBlock(in_ch=out_ch,out_ch=out_ch,stride=1,residual=self.residual))
        return nn.Sequential(*blocks)
\end{lstlisting}
Only the \emph{first} block of a stage gets the stage's stride and the channel change; the other $n-1$ are identity blocks. Stage 1 has stride 1 and $16\to16$, so all its blocks are identity blocks and there is no downsampling at the start. The forward pass is the straight composition
\[
  z = \text{Linear}\Big(\text{Flatten}\big(\text{GAP}(\text{stage}_3(\text{stage}_2(\text{stage}_1(\ReLU(\BN(\text{stem}(x)))))))\big)\Big).
\]
\end{repo}

\subsection{The shape flow, drawn}

\begin{center}
\begin{tikzpicture}[scale=0.88,ar/.style={-{Stealth[length=2mm]},thick}]
  % row 1
  \tensor{0}{0}{3.2}{3.2}{0.3}{cIn}{$32$}{$32$}{$3$}
  \node[font=\scriptsize,align=center] at (1.6,-0.95) {input\\\shp{N,3,32,32}};
  \draw[ar] (3.6,1.6) -- (4.6,1.6) node[midway,above,font=\scriptsize]{stem};
  \tensor{4.8}{0}{3.2}{3.2}{0.8}{c16}{$32$}{}{$16$}
  \node[font=\scriptsize,align=center] at (6.4,-0.95) {stem + BN + ReLU\\\shp{N,16,32,32}};
  \draw[ar] (8.6,1.6) -- (9.6,1.6) node[midway,above,font=\scriptsize]{$3$ blocks};
  \tensor{9.8}{0}{3.2}{3.2}{0.8}{c16}{$32$}{}{$16$}
  \node[font=\scriptsize,align=center] at (11.4,-0.95) {stage 1\\\shp{N,16,32,32}};
  \draw[ar] (13.6,1.6) -- (14.6,1.6) node[midway,above,font=\scriptsize]{$S{=}2$};
  \tensor{14.8}{0.8}{1.6}{1.6}{1.6}{c32}{$16$}{}{$32$}
  \node[font=\scriptsize,align=center] at (15.9,-0.95) {stage 2\\\shp{N,32,16,16}};
  % row 2
  \begin{scope}[shift={(0,-5.2)}]
  \draw[ar] (0,1.2) -- (1.2,1.2) node[midway,above,font=\scriptsize]{$S{=}2$};
  \tensor{1.4}{0.8}{0.8}{0.8}{3.2}{c64}{$8$}{}{$64$}
  \node[font=\scriptsize,align=center] at (2.6,-0.5) {stage 3\\\shp{N,64,8,8}};
  \draw[ar] (4.6,1.2) -- (5.8,1.2) node[midway,above,font=\scriptsize]{GAP};
  \tensor{6.0}{1.1}{0.12}{0.12}{3.2}{c64}{}{}{$64$}
  \node[font=\scriptsize,align=center] at (6.4,0.5) {\shp{N,64,1,1}};
  \draw[ar] (8.2,1.2) -- (9.2,1.2) node[midway,above,font=\scriptsize]{flatten};
  \matr{9.4}{1.05}{3.2}{0.3}{c64}{}{}{$p$: \shp{N,64}}
  \draw[ar] (12.8,1.2) -- (14.3,1.2) node[midway,above,font=\scriptsize]{$pW\T{+}b$};
  \matr{14.5}{1.05}{0.9}{0.3}{cG}{}{}{}
  \node[font=\scriptsize,align=center] at (14.95,0.5) {$Z$\\\shp{N,10}};
  \draw[ar] (15.6,1.2) -- (16.7,1.2) node[midway,above,font=\scriptsize]{$\argmax$};
  \node[draw,rounded corners,font=\scriptsize,fill=repo!15,align=center] at (17.4,1.2) {$\hat y$\\\shp{N}};
  \end{scope}
  \draw[ar,dashed,black!50] (15.9,-1.4) -- (15.9,-2.2) -- (0,-2.2) -- (0,-3.8);
\end{tikzpicture}
\end{center}

Read it as two trends running in opposite directions. The spatial size halves twice ($32\to16\to8$) and the depth doubles twice ($16\to32\to64$). The number of values per image shrinks from $16{,}384$ in stage 1 to $8{,}192$ in stage 2 and $4{,}096$ in stage 3: the network trades \emph{where} for \emph{what}.

\subsection{The full layer table}

\begin{center}
\small
\renewcommand{\arraystretch}{1.08}
\begin{longtable}{@{}rllrrr@{}}
\toprule
\# & layer (\code{state\_dict} prefix) & output shape & params & MACs / image & RF\\
\midrule
\endhead
 & input & \shp{N,3,32,32} & -- & -- & 1\\
1 & \code{stem} conv $3{\to}16$, $S1$ & \shp{N,16,32,32} & 432 & 442{,}368 & 3\\
 & \code{bn1}, ReLU & \shp{N,16,32,32} & 32 & -- & 3\\
\midrule
2--3 & \code{stage1.0} conv, BN, ReLU, conv, BN, $+x$, ReLU & \shp{N,16,32,32} & 4{,}672 & 4{,}718{,}592 & 7\\
4--5 & \code{stage1.1} & \shp{N,16,32,32} & 4{,}672 & 4{,}718{,}592 & 11\\
6--7 & \code{stage1.2} & \shp{N,16,32,32} & 4{,}672 & 4{,}718{,}592 & 15\\
\midrule
8--9 & \code{stage2.0} \textbf{$16{\to}32$, $S2$}, option-A shortcut & \shp{N,32,16,16} & 13{,}952 & 3{,}538{,}944 & 21\\
10--11 & \code{stage2.1} & \shp{N,32,16,16} & 18{,}560 & 4{,}718{,}592 & 29\\
12--13 & \code{stage2.2} & \shp{N,32,16,16} & 18{,}560 & 4{,}718{,}592 & 37\\
\midrule
14--15 & \code{stage3.0} \textbf{$32{\to}64$, $S2$}, option-A shortcut & \shp{N,64,8,8} & 55{,}552 & 3{,}538{,}944 & 49\\
16--17 & \code{stage3.1} & \shp{N,64,8,8} & 73{,}984 & 4{,}718{,}592 & 65\\
18--19 & \code{stage3.2} & \shp{N,64,8,8} & 73{,}984 & 4{,}718{,}592 & 81\\
\midrule
 & \code{avg\_pool} (GAP) & \shp{N,64,1,1} & 0 & -- & all\\
 & \code{flatten} & \shp{N,64} & 0 & -- & \\
20 & \code{linear} $64{\to}10$ & \shp{N,10} & 650 & 640 & \\
\midrule
 & \textbf{total} & & \textbf{269{,}722} & \textbf{40{,}551{,}040} & \\
\bottomrule
\end{longtable}
\end{center}

\begin{shapebox}[how each parameter count is obtained]
\begin{align*}
\text{stem} &: 16\cdot3\cdot3\cdot3 = 432, \qquad \text{BN}(16): 2\cdot16 = 32\\
\text{identity block, width } C &: 2\cdot(C\cdot C\cdot 9) + 2\cdot(2C) = 18C^2 + 4C\\
  &\;\Rightarrow\; C{=}16: 4{,}672,\quad C{=}32: 18{,}560,\quad C{=}64: 73{,}984\\
\text{downsampling block } C\to2C &: (2C\cdot C\cdot9) + (2C\cdot2C\cdot9) + 2\cdot(2\cdot2C) = 54C^2+8C\\
  &\;\Rightarrow\; 16{\to}32: 13{,}952,\quad 32{\to}64: 55{,}552\\
\text{linear} &: 64\cdot10 + 10 = 650\\
\text{total} &: 432 + 32 + 3(4{,}672) + (13{,}952 + 2\cdot18{,}560)\\
  &\quad + (55{,}552 + 2\cdot73{,}984) + 650 \;=\; \mathbf{269{,}722} .
\end{align*}
Per stage: stem $464$, stage 1 $14{,}016$, stage 2 $51{,}072$, stage 3 $203{,}520$, linear $650$. \textbf{Three quarters of all parameters live in stage 3}, because a conv's parameter count grows with $C^2$ while its spatial size does not enter. Of the 269{,}722 parameters, 267{,}696 are conv weights (19 convs), 1{,}376 are BN $\gamma/\beta$ (19 BN layers), and 650 are in the linear layer. The model also carries 1{,}376 BN running-statistic values (\code{running\_mean}, \code{running\_var}) and 19 \code{num\_batches\_tracked} counters as \emph{buffers}: they are saved in the checkpoint (116 \code{state\_dict} entries in total) but not trained by SGD.

The paper reports ``0.27M'' parameters for ResNet-20 and ``0.85M'' for ResNet-56; counting the model gives $269{,}722$ and $853{,}018$. The plain networks have exactly the same counts, because the shortcut is parameter-free.
\end{shapebox}

\paragraph{MACs.} Every $3{\times}3$ conv inside a stage costs $2{,}359{,}296$ MACs per image, whatever the stage (\cref{sec:im2col}). The two stride-2 convs cost half that ($1{,}179{,}648$), because they produce a quarter of the pixels with twice the filters of their input. The whole forward pass is $\approx 40.6$M MACs per image, $\approx 81$M FLOPs, or $\approx5.2$G MACs for a batch of 128. Backward costs roughly twice the forward.

\paragraph{Receptive field (RF).} The last column is the side length of the input patch that can influence one output pixel along the \emph{longest} (all-convolution) path. Each $3{\times}3$ conv adds $2\times$(current jump), and the jump doubles at each stride-2 conv: $3,5,\dots,15$ in stage 1 (jump 1), $+2$ then $+4$ per conv in stage 2, $+4$ then $+8$ in stage 3, ending at $81>32$. By stage 3 every unit can see the whole image plus padding, before GAP even averages over all positions. Shortcut paths have smaller receptive fields, so a ResNet behaves like an ensemble of paths of different lengths.

\subsection{Initialisation}
\label{sec:init}

\begin{lstlisting}[firstnumber=66]
    def _init_weights(self):
        for m in self.modules():
            if isinstance(m,nn.Conv2d):
                fan_in = m.weight[0].numel()
                std = (2/fan_in) ** 0.5
                with torch.no_grad():
                    m.weight.normal_(0,std)
\end{lstlisting}

\begin{idea}[He (Kaiming) initialisation, derived]
Take one output of a conv or linear layer, $y=\sum_{i=1}^{n}w_ix_i$, with $n=\text{fan\_in}$ inputs, weights i.i.d.\ with mean 0 and variance $\sigma^2$, independent of the inputs. Then
\[
  \operatorname{Var}(y) = n\,\sigma^2\,\mathbb E[x^2].
\]
If $x=\ReLU(z)$ with $z$ symmetric around 0, half of the mass is zeroed, so $\mathbb E[x^2]=\tfrac12\operatorname{Var}(z)$. Chaining layers,
\[
  \operatorname{Var}(y_l) = \tfrac12\,n_l\,\sigma_l^2\,\operatorname{Var}(y_{l-1}),
\]
which neither explodes nor vanishes with depth iff $\tfrac12 n_l\sigma_l^2=1$:
\[
  \boxed{\sigma_l = \sqrt{2/n_l}}, \qquad n_l = C_\text{in}\cdot K\cdot K .
\]
\end{idea}

\begin{repo}[what the code initialises, and what it leaves alone]
\code{m.weight[0]} is one filter, shape \shp{C_\text{in},3,3}, so \code{numel()} $=9C_\text{in}$ is exactly the fan-in.
\begin{center}
\small
\begin{tabular}{@{}lrrr@{}}
\toprule
conv & fan\_in $=9C_\text{in}$ & $\sigma=\sqrt{2/\text{fan\_in}}$ & std of an actual draw (seed 0)\\
\midrule
stem ($C_\text{in}{=}3$) & 27 & 0.2722 & 0.2708\\
$C_\text{in}{=}16$ (stage 1, first conv of stage 2) & 144 & 0.1179 & 0.1178 / 0.1172\\
$C_\text{in}{=}32$ (stage 2, first conv of stage 3) & 288 & 0.0833 & 0.0832\\
$C_\text{in}{=}64$ (stage 3) & 576 & 0.0589 & 0.0586\\
\bottomrule
\end{tabular}
\end{center}
\begin{itemize}[nosep]
  \item \textbf{Convolutions}: He-normal with fan-in, overwriting PyTorch's default \code{kaiming\_uniform\_(a=}$\sqrt5$\code{)}.
  \item \textbf{Linear layer}: \emph{not} touched, so it keeps PyTorch's default, $W,b\sim\mathcal U(-1/\sqrt{64},\,1/\sqrt{64})=\mathcal U(-0.125,0.125)$. The README states this explicitly.
  \item \textbf{BatchNorm}: PyTorch default, $\gamma=1$, $\beta=0$, running mean 0, running variance 1.
\end{itemize}
\end{repo}

\begin{warn}[fan-in vs.\ fan-out]
He et al.\ derive both a forward (fan-in) and a backward (fan-out, $n=9C_\text{out}$) condition. This repo uses fan-in. torchvision's reference ImageNet ResNet uses \code{mode="fan\_out"}. With BN after every conv the network is fairly insensitive to the choice, since BN re-normalises each channel on the first forward pass. The stem's input is not a ReLU output, so the factor 2 is not justified there; applying the same rule to every conv is the usual simplification.
\end{warn}

\subsection{Memory of one forward pass}

Per image, the outputs of the main ops are: stem/BN/ReLU $3\times16{,}384$; stage 1: 3 blocks $\times$ 7 tensors $\times$ $16{,}384$; stage 2 and 3 likewise with $8{,}192$ and $4{,}096$. That is $\approx 0.65$M activation values per image, about 2.6\,MB in \code{float32}, and on the order of 300\,MB for a batch of 128 (autograd saves most, but not exactly all, of them). Autograd keeps these for the backward pass. This, not the 1\,MB of parameters, is what fills GPU memory during training, and it is why \code{evaluate()} wraps the forward in \code{torch.no\_grad()}.
